% ── Section 3: Benchmark methodology ──────────────────────
\section{Benchmark Methodology}
\label{sec:methodology}

This section defines the benchmark used to compare reconstruction methods for the two sensor
targets. The benchmark asks how well different information sources can recover missing values in
a historical daily record: local target continuity, external physical covariates, gap-edge
context, and pretrained sequence modelling. Internal gaps are treated as historical imputation
problems, so observations on both sides of the missing interval may be used when available.

\subsection{Model ladder}
\label{sec:model_ladder}

Let $y_t$ denote the daily target on its modelling scale, and let
$\mathbf{x}_t$ denote the external covariate information available at day $t$. For an artificial
gap
\[
    G=\{t_0,\ldots,t_0+L-1\},
\]
the values $\mathbf y_G=\{y_t:t\in G\}$ are hidden and reconstructed as
$\hat{\mathbf y}_G=\{\hat y_t:t\in G\}$. All methods solve this same reconstruction problem, but
they differ in the information representation they use: simple temporal rules, target-only
continuity, external covariate features, explicit gap-edge context, or a pretrained sequence
model.

\begin{figure}[htbp]
  \centering
  \makebox[\textwidth][c]{%
    \hspace*{-1.5cm}%
    \resizebox{1.01\textwidth}{!}{% TikZ source: model-ladder hero figure with embedded generated subfigures.
% Included via:
%   \resizebox{\textwidth}{!}{% TikZ source: model-ladder hero figure with embedded generated subfigures.
% Included via:
%   \resizebox{\textwidth}{!}{% TikZ source: model-ladder hero figure with embedded generated subfigures.
% Included via:
%   \resizebox{\textwidth}{!}{\input{figures/fig_model_ladder_tikz.tex}}
%
% Requires graphicx in preamble.
% Expected external files (optional; fallback box shown if absent):
%   figures/model_ladder_external_tabular.png
%   figures/model_ladder_gap_edge_residual.png
%   figures/model_ladder_tsicl.png

\definecolor{vgTarget}{HTML}{2A7F8E}  % observed target (teal)
\definecolor{vgRecon}{HTML}{1B9E77}   % reconstruction / smoother output (green)
\definecolor{vgExt}{HTML}{3B6FA0}     % external predictors / features (blue)
\definecolor{vgEdge}{HTML}{D98A2B}    % gap-edge / residual correction (amber)
\definecolor{vgPrior}{HTML}{7E6CA8}   % TS-ICL / pretrained sequence model (violet)
\definecolor{vgHidden}{HTML}{C0504D}  % hidden gap (red)
\definecolor{vgInk}{HTML}{2B2B2B}     % text
\definecolor{vgFaint}{HTML}{9AA6B2}   % row tags / faint rules
\definecolor{vgGrid}{HTML}{D9DEE2}    % daily time grid
\definecolor{vgArrow}{HTML}{8A9299}   % flow arrows
\definecolor{vgPanel}{HTML}{F7F8F9}   % panel background
\definecolor{vgChip}{HTML}{ECEFF2}    % model chip fill

\definecolor{mBase}{HTML}{5F7F8A}     % simple baselines
\definecolor{mGP}{HTML}{1B9E77}       % Gaussian process
\definecolor{mExt}{HTML}{3B6FA0}      % external tabular
\definecolor{mGap}{HTML}{D98A2B}      % gap-edge residual
\definecolor{mTS}{HTML}{7E6CA8}       % TS-ICL
\definecolor{tsgreen}{HTML}{08A69B}
\definecolor{tsblue}{HTML}{004CF8}


\begin{tikzpicture}[x=0.38cm, y=0.38cm, font=\small,
    flow/.style={-{Latex[length=2mm]}, vgArrow, line width=0.9pt}]

  % ---------------------------------------------------------------------------
  % Global geometry
  % ---------------------------------------------------------------------------
  \def\PW{9.7}        % panel width
  \def\PH{24.9}       % panel height
  \def\DX{10.55}      % horizontal shift between panels

  \def\Ytitle{23.95}
  \def\YmodelA{20.60}
  \def\YmodelB{22.70}
  \def\YhowA{10.45}
  \def\YhowB{19.75}
  \def\YoutA{1.55}
  \def\YoutB{9.20}

  \def\Yeq{7.65}      % equation height inside output box
  \def\ob{2.95}       % baseline of output mini-plots

  % ---------------------------------------------------------------------------
  % Helpers
  % ---------------------------------------------------------------------------

  \newcommand{\panelbg}[2]{%
    \fill[vgPanel, rounded corners=3pt] (-0.25,0.75) rectangle (\PW,\PH);
    \draw[#2, opacity=0.45, line width=0.65pt, rounded corners=3pt]
      (-0.25,0.75) rectangle (\PW,\PH);
    \node[font=\small\bfseries, text=vgInk, align=center] at (4.72,\Ytitle) {#1};
  }

  \newcommand{\modelchip}[1]{%
    \fill[vgChip, rounded corners=3pt] (1.05,\YmodelA) rectangle (8.40,\YmodelB);
    \draw[#1, opacity=0.40, line width=0.55pt, rounded corners=3pt]
      (1.05,\YmodelA) rectangle (8.40,\YmodelB);
  }

  \newcommand{\howbox}[1]{%
    \fill[#1, opacity=0.040, rounded corners=3pt] (0.55,\YhowA) rectangle (8.95,\YhowB);
    \draw[#1, opacity=0.75, line width=0.7pt, rounded corners=3pt]
      (0.55,\YhowA) rectangle (8.95,\YhowB);
  }

  \newcommand{\outbox}[1]{%
    \fill[#1, opacity=0.040, rounded corners=3pt] (0.55,\YoutA) rectangle (8.95,\YoutB);
    \draw[#1, opacity=0.75, line width=0.7pt, rounded corners=3pt]
      (0.55,\YoutA) rectangle (8.95,\YoutB);
  }

  \newcommand{\imageorfallback}[3]{%
    \IfFileExists{#1}{%
      \includegraphics[width=#2]{#1}%
    }{%
      \fbox{\parbox[c][2.8cm][c]{#2}{\centering\scriptsize #3}}%
    }%
  }

  % output grid with observed context and hidden interval
  \newcommand{\miniout}[1]{%
    \foreach \s in {1,...,9}{
      \draw[vgGrid,line width=0.28pt](0.2+\s*0.9,#1)--(0.2+\s*0.9,#1+3.3);
    }
    \fill[vgHidden,opacity=0.12](0.2+3.5*0.9,#1)rectangle(0.2+6.5*0.9,#1+3.3);
    \draw[vgTarget,line width=1.0pt](1.1,#1+1.35)--(2.0,#1+1.00)--(2.9,#1+1.20);
    \draw[vgTarget,line width=1.0pt](6.5,#1+2.50)--(7.4,#1+2.85)--(8.3,#1+2.50);
    \foreach \s/\y in {1/1.35,2/1.00,3/1.20,7/2.50,8/2.85,9/2.50}{
      \fill[vgTarget](0.2+\s*0.9,#1+\y)circle(0.10);
    }
  }

  % compact target-gap sketch used in "how" boxes
  \newcommand{\howgap}[1]{%
    \foreach \s in {1,...,9}{
      \draw[vgGrid,line width=0.25pt](0.2+\s*0.9,#1)--(0.2+\s*0.9,#1+3.3);
    }
    \fill[vgHidden,opacity=0.12](0.2+3.5*0.9,#1)rectangle(0.2+6.5*0.9,#1+3.3);
    \draw[vgTarget,line width=0.95pt](1.1,#1+1.35)--(2.0,#1+1.00)--(2.9,#1+1.20);
    \draw[vgTarget,line width=0.95pt](6.5,#1+2.50)--(7.4,#1+2.85)--(8.3,#1+2.50);
    \foreach \s/\y in {1/1.35,2/1.00,3/1.20,7/2.50,8/2.85,9/2.50}{
      \fill[vgTarget](0.2+\s*0.9,#1+\y)circle(0.09);
    }
  }

  % left row labels
  \node[rotate=90,font=\footnotesize\bfseries,text=vgFaint] at (-2.25,22.10) {METHOD};
  \node[rotate=90,font=\footnotesize\bfseries,text=vgFaint] at (-2.25,15.10) {MECHANISM};
  \node[rotate=90,font=\footnotesize\bfseries,text=vgFaint] at (-2.25,5.35) {OUTPUT};

  % ---------------------------------------------------------------------------
  % Panel 1: Simple baselines
  % ---------------------------------------------------------------------------
  \begin{scope}[shift={(0,0)}]
    \panelbg{Simple baselines}{mBase}
    \modelchip{mBase}
    \node[font=\tiny, text=vgInk, align=center, text width=3.2cm] at (4.72,21.65)
      {climatology\\persistence\\linear interpolation};

    \howbox{mBase}
    % three rule sketches with equations, not names
    \def\yshift{-0.9}
\foreach \yy/\eqn in {
  17.90/{{$\hat y_t = y_{t^-}$}},
  15.90/{{$\hat y_t=\tilde y_t^{\mathrm{lin}}$}},
  13.90/{{$\hat y_t = \mu_{\mathrm{season}}$}}
}{
  \fill[vgHidden,opacity=0.12](1.35,\yy-0.55+\yshift) rectangle (2.30,\yy+0.55+\yshift);
  \fill[vgTarget](0.95,\yy-0.25+\yshift) circle(0.09);
  \fill[vgTarget](2.80,\yy+0.30+\yshift) circle(0.09);
  \node[font=\scriptsize,text=vgInk,anchor=west] at (3.60,\yy+\yshift) {\eqn};
}
    % persistence
    \draw[mBase,line width=1.35pt] (0.95,16.75)--(2.80,16.75);
    % interpolation
    \draw[mBase,line width=1.35pt] (0.95,14.75)--(2.80,15.30);
    % climatology
    \draw[mBase,line width=1.35pt] (0.95,13)--(2.80,13);

    \draw[flow] (4.75,20.35)--(4.75,19.95);
    \draw[flow] (4.75,10.25)--(4.75,9.55);

    \outbox{mBase}
    \node[font=\scriptsize,text=vgInk,align=center] at (4.75,\Yeq)
      {$\hat y_t=\tilde y_t^{\mathrm{lin}}$};
    \miniout{\ob}
    % only interpolation shown in output
    \draw[mBase,line width=1.6pt](2.9,\ob+1.2)--(6.5,\ob+2.5);
    \foreach \s in {4,5,6}{
      \fill[mBase](0.2+\s*0.9,{\ob+1.2+(\s-3)*0.325})circle(0.10);
    }
  \end{scope}

  % ---------------------------------------------------------------------------
  % Panel 2: Target-only smoother
  % ---------------------------------------------------------------------------
  \begin{scope}[shift={(\DX,0)}]
    \panelbg{\scalebox{0.93}{Target-only smoother}}{mGP}
    \modelchip{mGP}
    \node[font=\scriptsize, text=vgInk, align=center, text width=3.0cm] at (4.72,21.65)
      {Gaussian process};

    \howbox{mGP}

\begin{scope}[yshift=-0.35cm]
\howgap{13.85}

% Prior mean m(t): faint dashed reference trend
\draw[vgFaint,line width=0.75pt,dash pattern=on 2pt off 1.5pt]
  plot[smooth] coordinates{
    (1.1,15.20)
    (2.9,15.10)
    (4.7,15.70)
    (6.5,16.35)
    (8.3,16.25)
  };
\node[font=\scriptsize,text=vgFaint,anchor=west] at (6.75,14.75) {$m(t)$};

% Posterior uncertainty induced by the covariance kernel inside/near the gap
\fill[mGP,opacity=0.13]
  plot[smooth] coordinates{
    (2.9,15.32)
    (3.6,15.55)
    (4.7,16.25)
    (5.8,16.55)
    (6.5,16.58)
  }
  -- plot[smooth] coordinates{
    (6.5,16.12)
    (5.8,15.55)
    (4.7,15.05)
    (3.6,14.92)
    (2.9,14.92)
  }
  -- cycle;

% Posterior mean reconstruction
\draw[mGP,line width=1.65pt]
  plot[smooth] coordinates{
    (2.9,15.05)
    (3.6,15.18)
    (4.7,15.75)
    (5.8,16.12)
    (6.5,16.35)
  };

% Small covariance cue
\draw[mGP,line width=0.55pt,opacity=0.85]
  (4.15,16.15) to[bend left=30] (5.25,16.15);
\draw[mGP,line width=0.55pt,opacity=0.65]
  (3.80,16.42) to[bend left=35] (5.60,16.42);
\node[font=\scriptsize,text=mGP] at (4.70,17.80) {$k(t,t')$};
\end{scope}

\draw[flow] (4.75,20.35)--(4.75,19.95);
\draw[flow] (4.75,10.25)--(4.75,9.55);

    \outbox{mGP}
    \node[font=\scriptsize,text=vgInk,align=center] at (4.75,\Yeq)
      {$\hat y(t)\sim\mathcal{GP}(m,k)$};
    \miniout{\ob}
    \fill[mGP,opacity=0.14]
      plot[smooth]coordinates{(2.9,\ob+1.25)(3.8,\ob+1.20)(4.7,\ob+2.75)(5.6,\ob+2.25)(6.5,\ob+2.55)}
      --plot[smooth]coordinates{(6.5,\ob+2.45)(5.6,\ob+1.55)(4.7,\ob+2.00)(3.8,\ob+0.60)(2.9,\ob+1.15)}
      --cycle;
    \draw[mGP,line width=1.6pt]
      plot[smooth]coordinates{(2.9,\ob+1.20)(3.8,\ob+0.90)(4.7,\ob+2.40)(5.6,\ob+1.90)(6.5,\ob+2.50)};
  \end{scope}

  % ---------------------------------------------------------------------------
  % Panel 3: External tabular
  % ---------------------------------------------------------------------------
  \begin{scope}[shift={({2*\DX},0)}]
    \panelbg{External tabular}{mExt}
    \modelchip{mExt}
    \node[font=\scriptsize, text=vgInk, align=center, text width=3.0cm] at (4.72,21.65)
      {HGB, ExtraTrees};

    \howbox{mExt}
    \node[inner sep=0pt] at (4.75,15.40)
      {\imageorfallback{figures/model_ladder_external_tabular.png}{3.35cm}{external tabular subfigure}};
    % symbol overlay
    \node[font=\scriptsize,text=mExt] at (6.40,12.6) {$\phi_t$};

    \draw[flow] (4.75,20.35)--(4.75,19.95);
    \draw[flow] (4.75,10.25)--(4.75,9.55);

    \outbox{mExt}
    \node[font=\scriptsize,text=vgInk,align=center] at (4.75,\Yeq)
      {$\hat y_t = f(\phi_t)$};
    \miniout{\ob}
    \foreach \s/\y in {4/1.35,5/1.90,6/1.55}{
      \fill[mExt](0.2+\s*0.9,\ob+\y)circle(0.15);
    }
  \end{scope}

  % ---------------------------------------------------------------------------
  % Panel 4: Gap-edge residual
  % ---------------------------------------------------------------------------
  \begin{scope}[shift={({3*\DX},0)}]
    \panelbg{Gap-edge residual}{mGap}
    \modelchip{mGap}
    \node[font=\scriptsize, text=vgInk, align=center, text width=3.0cm] at (4.72,21.65)
      {HGB, ExtraTrees};

    \howbox{mGap}
    \node[inner sep=0pt] at (4.75,15.40)
      {\imageorfallback{figures/model_ladder_gap_edge_residual.png}{3.30cm}{gap-edge residual subfigure}};
    % symbol overlays
    \node[font=\scriptsize,text=mExt] at (6.72,12.9) {$\phi_t$};
    \node[font=\scriptsize,text=mGap] at (1.7,12) {$\psi_G$};

    \draw[flow] (4.75,20.35)--(4.75,19.95);
    \draw[flow] (4.75,10.25)--(4.75,9.55);

    \outbox{mGap}
    \node[font=\scriptsize,text=vgInk,align=center] at (4.75,\Yeq)
      {$\hat y_t=\tilde y_t^{\mathrm{lin}} + g(\phi_t,\psi_G)$};
    \miniout{\ob}
    % interpolation baseline
    \draw[vgArrow,line width=1.1pt,dash pattern=on 1.7pt off 1.4pt](2.9,\ob+1.2)--(6.5,\ob+2.5);
    % residual bars only; no connecting curve
    \foreach \s/\by/\cy in {4/1.525/0.95,5/1.85/2.42,6/2.175/1.92}{
      \draw[mGap,line width=1.35pt](0.2+\s*0.9,\ob+\by)--(0.2+\s*0.9,\ob+\cy);
      \fill[mGap](0.2+\s*0.9,\ob+\cy)circle(0.11);
    }
  \end{scope}

  % ---------------------------------------------------------------------------
  % Panel 5: TS-ICL
  % ---------------------------------------------------------------------------
  \begin{scope}[shift={({4*\DX},0)}]
    \panelbg{TS-ICL}{mTS}
    \modelchip{mTS}
    \node[font=\scriptsize, text=vgInk, align=center, text width=3.0cm] at (4.72,21.65)
      {pretrained\\zero-shot};

    \howbox{mTS}
    \node[inner sep=0pt] at (4.75,14.9)
      {\imageorfallback{figures/model_ladder_tsicl.png}{2.75cm}{TS-ICL subfigure}};
    % symbol overlays
    \node[font=\scriptsize,text=tsgreen,anchor=west] at (1.3,19) {$ \mathbf y_{\mathrm{ctxt}}$};
    \node[font=\scriptsize,text=tsblue,anchor=west] at (3.1,16) {$\mathbf X_{\mathrm{covar}}$};

    \draw[flow] (4.75,20.35)--(4.75,19.95);
    \draw[flow] (4.75,10.25)--(4.75,9.55);

    \outbox{mTS}
    \node[font=\scriptsize,text=vgInk,align=center] at (4.75,\Yeq)
      {$ \hat y(t)= F(\mathbf y_{\mathrm{ctxt}}, \mathbf X_{\mathrm{covar}})$};
    \miniout{\ob}
    \fill[mTS,opacity=0.13]
      plot[smooth]coordinates{(2.9,\ob+1.25)(3.8,\ob+1.25)(4.7,\ob+2.55)(5.6,\ob+2.20)(6.5,\ob+2.55)}
      --plot[smooth]coordinates{(6.5,\ob+2.45)(5.6,\ob+1.70)(4.7,\ob+1.85)(3.8,\ob+0.85)(2.9,\ob+1.15)}
      --cycle;
    \draw[mTS,line width=1.6pt]
      plot[smooth]coordinates{(2.9,\ob+1.20)(3.8,\ob+1.05)(4.7,\ob+2.20)(5.6,\ob+1.95)(6.5,\ob+2.50)};
  \end{scope}

  % ---------------------------------------------------------------------------
  % Top arrow
  % ---------------------------------------------------------------------------
  \draw[-{Latex[length=3mm]},vgFaint,line width=1.35pt](0.2,25.45)--(51.35,25.45);
  \node[font=\small\itshape,text=vgInk] at (25.75,26.20)
    {increasing information and modelling structure};

\end{tikzpicture}}
%
% Requires graphicx in preamble.
% Expected external files (optional; fallback box shown if absent):
%   figures/model_ladder_external_tabular.png
%   figures/model_ladder_gap_edge_residual.png
%   figures/model_ladder_tsicl.png

\definecolor{vgTarget}{HTML}{2A7F8E}  % observed target (teal)
\definecolor{vgRecon}{HTML}{1B9E77}   % reconstruction / smoother output (green)
\definecolor{vgExt}{HTML}{3B6FA0}     % external predictors / features (blue)
\definecolor{vgEdge}{HTML}{D98A2B}    % gap-edge / residual correction (amber)
\definecolor{vgPrior}{HTML}{7E6CA8}   % TS-ICL / pretrained sequence model (violet)
\definecolor{vgHidden}{HTML}{C0504D}  % hidden gap (red)
\definecolor{vgInk}{HTML}{2B2B2B}     % text
\definecolor{vgFaint}{HTML}{9AA6B2}   % row tags / faint rules
\definecolor{vgGrid}{HTML}{D9DEE2}    % daily time grid
\definecolor{vgArrow}{HTML}{8A9299}   % flow arrows
\definecolor{vgPanel}{HTML}{F7F8F9}   % panel background
\definecolor{vgChip}{HTML}{ECEFF2}    % model chip fill

\definecolor{mBase}{HTML}{5F7F8A}     % simple baselines
\definecolor{mGP}{HTML}{1B9E77}       % Gaussian process
\definecolor{mExt}{HTML}{3B6FA0}      % external tabular
\definecolor{mGap}{HTML}{D98A2B}      % gap-edge residual
\definecolor{mTS}{HTML}{7E6CA8}       % TS-ICL
\definecolor{tsgreen}{HTML}{08A69B}
\definecolor{tsblue}{HTML}{004CF8}


\begin{tikzpicture}[x=0.38cm, y=0.38cm, font=\small,
    flow/.style={-{Latex[length=2mm]}, vgArrow, line width=0.9pt}]

  % ---------------------------------------------------------------------------
  % Global geometry
  % ---------------------------------------------------------------------------
  \def\PW{9.7}        % panel width
  \def\PH{24.9}       % panel height
  \def\DX{10.55}      % horizontal shift between panels

  \def\Ytitle{23.95}
  \def\YmodelA{20.60}
  \def\YmodelB{22.70}
  \def\YhowA{10.45}
  \def\YhowB{19.75}
  \def\YoutA{1.55}
  \def\YoutB{9.20}

  \def\Yeq{7.65}      % equation height inside output box
  \def\ob{2.95}       % baseline of output mini-plots

  % ---------------------------------------------------------------------------
  % Helpers
  % ---------------------------------------------------------------------------

  \newcommand{\panelbg}[2]{%
    \fill[vgPanel, rounded corners=3pt] (-0.25,0.75) rectangle (\PW,\PH);
    \draw[#2, opacity=0.45, line width=0.65pt, rounded corners=3pt]
      (-0.25,0.75) rectangle (\PW,\PH);
    \node[font=\small\bfseries, text=vgInk, align=center] at (4.72,\Ytitle) {#1};
  }

  \newcommand{\modelchip}[1]{%
    \fill[vgChip, rounded corners=3pt] (1.05,\YmodelA) rectangle (8.40,\YmodelB);
    \draw[#1, opacity=0.40, line width=0.55pt, rounded corners=3pt]
      (1.05,\YmodelA) rectangle (8.40,\YmodelB);
  }

  \newcommand{\howbox}[1]{%
    \fill[#1, opacity=0.040, rounded corners=3pt] (0.55,\YhowA) rectangle (8.95,\YhowB);
    \draw[#1, opacity=0.75, line width=0.7pt, rounded corners=3pt]
      (0.55,\YhowA) rectangle (8.95,\YhowB);
  }

  \newcommand{\outbox}[1]{%
    \fill[#1, opacity=0.040, rounded corners=3pt] (0.55,\YoutA) rectangle (8.95,\YoutB);
    \draw[#1, opacity=0.75, line width=0.7pt, rounded corners=3pt]
      (0.55,\YoutA) rectangle (8.95,\YoutB);
  }

  \newcommand{\imageorfallback}[3]{%
    \IfFileExists{#1}{%
      \includegraphics[width=#2]{#1}%
    }{%
      \fbox{\parbox[c][2.8cm][c]{#2}{\centering\scriptsize #3}}%
    }%
  }

  % output grid with observed context and hidden interval
  \newcommand{\miniout}[1]{%
    \foreach \s in {1,...,9}{
      \draw[vgGrid,line width=0.28pt](0.2+\s*0.9,#1)--(0.2+\s*0.9,#1+3.3);
    }
    \fill[vgHidden,opacity=0.12](0.2+3.5*0.9,#1)rectangle(0.2+6.5*0.9,#1+3.3);
    \draw[vgTarget,line width=1.0pt](1.1,#1+1.35)--(2.0,#1+1.00)--(2.9,#1+1.20);
    \draw[vgTarget,line width=1.0pt](6.5,#1+2.50)--(7.4,#1+2.85)--(8.3,#1+2.50);
    \foreach \s/\y in {1/1.35,2/1.00,3/1.20,7/2.50,8/2.85,9/2.50}{
      \fill[vgTarget](0.2+\s*0.9,#1+\y)circle(0.10);
    }
  }

  % compact target-gap sketch used in "how" boxes
  \newcommand{\howgap}[1]{%
    \foreach \s in {1,...,9}{
      \draw[vgGrid,line width=0.25pt](0.2+\s*0.9,#1)--(0.2+\s*0.9,#1+3.3);
    }
    \fill[vgHidden,opacity=0.12](0.2+3.5*0.9,#1)rectangle(0.2+6.5*0.9,#1+3.3);
    \draw[vgTarget,line width=0.95pt](1.1,#1+1.35)--(2.0,#1+1.00)--(2.9,#1+1.20);
    \draw[vgTarget,line width=0.95pt](6.5,#1+2.50)--(7.4,#1+2.85)--(8.3,#1+2.50);
    \foreach \s/\y in {1/1.35,2/1.00,3/1.20,7/2.50,8/2.85,9/2.50}{
      \fill[vgTarget](0.2+\s*0.9,#1+\y)circle(0.09);
    }
  }

  % left row labels
  \node[rotate=90,font=\footnotesize\bfseries,text=vgFaint] at (-2.25,22.10) {METHOD};
  \node[rotate=90,font=\footnotesize\bfseries,text=vgFaint] at (-2.25,15.10) {MECHANISM};
  \node[rotate=90,font=\footnotesize\bfseries,text=vgFaint] at (-2.25,5.35) {OUTPUT};

  % ---------------------------------------------------------------------------
  % Panel 1: Simple baselines
  % ---------------------------------------------------------------------------
  \begin{scope}[shift={(0,0)}]
    \panelbg{Simple baselines}{mBase}
    \modelchip{mBase}
    \node[font=\tiny, text=vgInk, align=center, text width=3.2cm] at (4.72,21.65)
      {climatology\\persistence\\linear interpolation};

    \howbox{mBase}
    % three rule sketches with equations, not names
    \def\yshift{-0.9}
\foreach \yy/\eqn in {
  17.90/{{$\hat y_t = y_{t^-}$}},
  15.90/{{$\hat y_t=\tilde y_t^{\mathrm{lin}}$}},
  13.90/{{$\hat y_t = \mu_{\mathrm{season}}$}}
}{
  \fill[vgHidden,opacity=0.12](1.35,\yy-0.55+\yshift) rectangle (2.30,\yy+0.55+\yshift);
  \fill[vgTarget](0.95,\yy-0.25+\yshift) circle(0.09);
  \fill[vgTarget](2.80,\yy+0.30+\yshift) circle(0.09);
  \node[font=\scriptsize,text=vgInk,anchor=west] at (3.60,\yy+\yshift) {\eqn};
}
    % persistence
    \draw[mBase,line width=1.35pt] (0.95,16.75)--(2.80,16.75);
    % interpolation
    \draw[mBase,line width=1.35pt] (0.95,14.75)--(2.80,15.30);
    % climatology
    \draw[mBase,line width=1.35pt] (0.95,13)--(2.80,13);

    \draw[flow] (4.75,20.35)--(4.75,19.95);
    \draw[flow] (4.75,10.25)--(4.75,9.55);

    \outbox{mBase}
    \node[font=\scriptsize,text=vgInk,align=center] at (4.75,\Yeq)
      {$\hat y_t=\tilde y_t^{\mathrm{lin}}$};
    \miniout{\ob}
    % only interpolation shown in output
    \draw[mBase,line width=1.6pt](2.9,\ob+1.2)--(6.5,\ob+2.5);
    \foreach \s in {4,5,6}{
      \fill[mBase](0.2+\s*0.9,{\ob+1.2+(\s-3)*0.325})circle(0.10);
    }
  \end{scope}

  % ---------------------------------------------------------------------------
  % Panel 2: Target-only smoother
  % ---------------------------------------------------------------------------
  \begin{scope}[shift={(\DX,0)}]
    \panelbg{\scalebox{0.93}{Target-only smoother}}{mGP}
    \modelchip{mGP}
    \node[font=\scriptsize, text=vgInk, align=center, text width=3.0cm] at (4.72,21.65)
      {Gaussian process};

    \howbox{mGP}

\begin{scope}[yshift=-0.35cm]
\howgap{13.85}

% Prior mean m(t): faint dashed reference trend
\draw[vgFaint,line width=0.75pt,dash pattern=on 2pt off 1.5pt]
  plot[smooth] coordinates{
    (1.1,15.20)
    (2.9,15.10)
    (4.7,15.70)
    (6.5,16.35)
    (8.3,16.25)
  };
\node[font=\scriptsize,text=vgFaint,anchor=west] at (6.75,14.75) {$m(t)$};

% Posterior uncertainty induced by the covariance kernel inside/near the gap
\fill[mGP,opacity=0.13]
  plot[smooth] coordinates{
    (2.9,15.32)
    (3.6,15.55)
    (4.7,16.25)
    (5.8,16.55)
    (6.5,16.58)
  }
  -- plot[smooth] coordinates{
    (6.5,16.12)
    (5.8,15.55)
    (4.7,15.05)
    (3.6,14.92)
    (2.9,14.92)
  }
  -- cycle;

% Posterior mean reconstruction
\draw[mGP,line width=1.65pt]
  plot[smooth] coordinates{
    (2.9,15.05)
    (3.6,15.18)
    (4.7,15.75)
    (5.8,16.12)
    (6.5,16.35)
  };

% Small covariance cue
\draw[mGP,line width=0.55pt,opacity=0.85]
  (4.15,16.15) to[bend left=30] (5.25,16.15);
\draw[mGP,line width=0.55pt,opacity=0.65]
  (3.80,16.42) to[bend left=35] (5.60,16.42);
\node[font=\scriptsize,text=mGP] at (4.70,17.80) {$k(t,t')$};
\end{scope}

\draw[flow] (4.75,20.35)--(4.75,19.95);
\draw[flow] (4.75,10.25)--(4.75,9.55);

    \outbox{mGP}
    \node[font=\scriptsize,text=vgInk,align=center] at (4.75,\Yeq)
      {$\hat y(t)\sim\mathcal{GP}(m,k)$};
    \miniout{\ob}
    \fill[mGP,opacity=0.14]
      plot[smooth]coordinates{(2.9,\ob+1.25)(3.8,\ob+1.20)(4.7,\ob+2.75)(5.6,\ob+2.25)(6.5,\ob+2.55)}
      --plot[smooth]coordinates{(6.5,\ob+2.45)(5.6,\ob+1.55)(4.7,\ob+2.00)(3.8,\ob+0.60)(2.9,\ob+1.15)}
      --cycle;
    \draw[mGP,line width=1.6pt]
      plot[smooth]coordinates{(2.9,\ob+1.20)(3.8,\ob+0.90)(4.7,\ob+2.40)(5.6,\ob+1.90)(6.5,\ob+2.50)};
  \end{scope}

  % ---------------------------------------------------------------------------
  % Panel 3: External tabular
  % ---------------------------------------------------------------------------
  \begin{scope}[shift={({2*\DX},0)}]
    \panelbg{External tabular}{mExt}
    \modelchip{mExt}
    \node[font=\scriptsize, text=vgInk, align=center, text width=3.0cm] at (4.72,21.65)
      {HGB, ExtraTrees};

    \howbox{mExt}
    \node[inner sep=0pt] at (4.75,15.40)
      {\imageorfallback{figures/model_ladder_external_tabular.png}{3.35cm}{external tabular subfigure}};
    % symbol overlay
    \node[font=\scriptsize,text=mExt] at (6.40,12.6) {$\phi_t$};

    \draw[flow] (4.75,20.35)--(4.75,19.95);
    \draw[flow] (4.75,10.25)--(4.75,9.55);

    \outbox{mExt}
    \node[font=\scriptsize,text=vgInk,align=center] at (4.75,\Yeq)
      {$\hat y_t = f(\phi_t)$};
    \miniout{\ob}
    \foreach \s/\y in {4/1.35,5/1.90,6/1.55}{
      \fill[mExt](0.2+\s*0.9,\ob+\y)circle(0.15);
    }
  \end{scope}

  % ---------------------------------------------------------------------------
  % Panel 4: Gap-edge residual
  % ---------------------------------------------------------------------------
  \begin{scope}[shift={({3*\DX},0)}]
    \panelbg{Gap-edge residual}{mGap}
    \modelchip{mGap}
    \node[font=\scriptsize, text=vgInk, align=center, text width=3.0cm] at (4.72,21.65)
      {HGB, ExtraTrees};

    \howbox{mGap}
    \node[inner sep=0pt] at (4.75,15.40)
      {\imageorfallback{figures/model_ladder_gap_edge_residual.png}{3.30cm}{gap-edge residual subfigure}};
    % symbol overlays
    \node[font=\scriptsize,text=mExt] at (6.72,12.9) {$\phi_t$};
    \node[font=\scriptsize,text=mGap] at (1.7,12) {$\psi_G$};

    \draw[flow] (4.75,20.35)--(4.75,19.95);
    \draw[flow] (4.75,10.25)--(4.75,9.55);

    \outbox{mGap}
    \node[font=\scriptsize,text=vgInk,align=center] at (4.75,\Yeq)
      {$\hat y_t=\tilde y_t^{\mathrm{lin}} + g(\phi_t,\psi_G)$};
    \miniout{\ob}
    % interpolation baseline
    \draw[vgArrow,line width=1.1pt,dash pattern=on 1.7pt off 1.4pt](2.9,\ob+1.2)--(6.5,\ob+2.5);
    % residual bars only; no connecting curve
    \foreach \s/\by/\cy in {4/1.525/0.95,5/1.85/2.42,6/2.175/1.92}{
      \draw[mGap,line width=1.35pt](0.2+\s*0.9,\ob+\by)--(0.2+\s*0.9,\ob+\cy);
      \fill[mGap](0.2+\s*0.9,\ob+\cy)circle(0.11);
    }
  \end{scope}

  % ---------------------------------------------------------------------------
  % Panel 5: TS-ICL
  % ---------------------------------------------------------------------------
  \begin{scope}[shift={({4*\DX},0)}]
    \panelbg{TS-ICL}{mTS}
    \modelchip{mTS}
    \node[font=\scriptsize, text=vgInk, align=center, text width=3.0cm] at (4.72,21.65)
      {pretrained\\zero-shot};

    \howbox{mTS}
    \node[inner sep=0pt] at (4.75,14.9)
      {\imageorfallback{figures/model_ladder_tsicl.png}{2.75cm}{TS-ICL subfigure}};
    % symbol overlays
    \node[font=\scriptsize,text=tsgreen,anchor=west] at (1.3,19) {$ \mathbf y_{\mathrm{ctxt}}$};
    \node[font=\scriptsize,text=tsblue,anchor=west] at (3.1,16) {$\mathbf X_{\mathrm{covar}}$};

    \draw[flow] (4.75,20.35)--(4.75,19.95);
    \draw[flow] (4.75,10.25)--(4.75,9.55);

    \outbox{mTS}
    \node[font=\scriptsize,text=vgInk,align=center] at (4.75,\Yeq)
      {$ \hat y(t)= F(\mathbf y_{\mathrm{ctxt}}, \mathbf X_{\mathrm{covar}})$};
    \miniout{\ob}
    \fill[mTS,opacity=0.13]
      plot[smooth]coordinates{(2.9,\ob+1.25)(3.8,\ob+1.25)(4.7,\ob+2.55)(5.6,\ob+2.20)(6.5,\ob+2.55)}
      --plot[smooth]coordinates{(6.5,\ob+2.45)(5.6,\ob+1.70)(4.7,\ob+1.85)(3.8,\ob+0.85)(2.9,\ob+1.15)}
      --cycle;
    \draw[mTS,line width=1.6pt]
      plot[smooth]coordinates{(2.9,\ob+1.20)(3.8,\ob+1.05)(4.7,\ob+2.20)(5.6,\ob+1.95)(6.5,\ob+2.50)};
  \end{scope}

  % ---------------------------------------------------------------------------
  % Top arrow
  % ---------------------------------------------------------------------------
  \draw[-{Latex[length=3mm]},vgFaint,line width=1.35pt](0.2,25.45)--(51.35,25.45);
  \node[font=\small\itshape,text=vgInk] at (25.75,26.20)
    {increasing information and modelling structure};

\end{tikzpicture}}
%
% Requires graphicx in preamble.
% Expected external files (optional; fallback box shown if absent):
%   figures/model_ladder_external_tabular.png
%   figures/model_ladder_gap_edge_residual.png
%   figures/model_ladder_tsicl.png

\definecolor{vgTarget}{HTML}{2A7F8E}  % observed target (teal)
\definecolor{vgRecon}{HTML}{1B9E77}   % reconstruction / smoother output (green)
\definecolor{vgExt}{HTML}{3B6FA0}     % external predictors / features (blue)
\definecolor{vgEdge}{HTML}{D98A2B}    % gap-edge / residual correction (amber)
\definecolor{vgPrior}{HTML}{7E6CA8}   % TS-ICL / pretrained sequence model (violet)
\definecolor{vgHidden}{HTML}{C0504D}  % hidden gap (red)
\definecolor{vgInk}{HTML}{2B2B2B}     % text
\definecolor{vgFaint}{HTML}{9AA6B2}   % row tags / faint rules
\definecolor{vgGrid}{HTML}{D9DEE2}    % daily time grid
\definecolor{vgArrow}{HTML}{8A9299}   % flow arrows
\definecolor{vgPanel}{HTML}{F7F8F9}   % panel background
\definecolor{vgChip}{HTML}{ECEFF2}    % model chip fill

\definecolor{mBase}{HTML}{5F7F8A}     % simple baselines
\definecolor{mGP}{HTML}{1B9E77}       % Gaussian process
\definecolor{mExt}{HTML}{3B6FA0}      % external tabular
\definecolor{mGap}{HTML}{D98A2B}      % gap-edge residual
\definecolor{mTS}{HTML}{7E6CA8}       % TS-ICL
\definecolor{tsgreen}{HTML}{08A69B}
\definecolor{tsblue}{HTML}{004CF8}


\begin{tikzpicture}[x=0.38cm, y=0.38cm, font=\small,
    flow/.style={-{Latex[length=2mm]}, vgArrow, line width=0.9pt}]

  % ---------------------------------------------------------------------------
  % Global geometry
  % ---------------------------------------------------------------------------
  \def\PW{9.7}        % panel width
  \def\PH{24.9}       % panel height
  \def\DX{10.55}      % horizontal shift between panels

  \def\Ytitle{23.95}
  \def\YmodelA{20.60}
  \def\YmodelB{22.70}
  \def\YhowA{10.45}
  \def\YhowB{19.75}
  \def\YoutA{1.55}
  \def\YoutB{9.20}

  \def\Yeq{7.65}      % equation height inside output box
  \def\ob{2.95}       % baseline of output mini-plots

  % ---------------------------------------------------------------------------
  % Helpers
  % ---------------------------------------------------------------------------

  \newcommand{\panelbg}[2]{%
    \fill[vgPanel, rounded corners=3pt] (-0.25,0.75) rectangle (\PW,\PH);
    \draw[#2, opacity=0.45, line width=0.65pt, rounded corners=3pt]
      (-0.25,0.75) rectangle (\PW,\PH);
    \node[font=\small\bfseries, text=vgInk, align=center] at (4.72,\Ytitle) {#1};
  }

  \newcommand{\modelchip}[1]{%
    \fill[vgChip, rounded corners=3pt] (1.05,\YmodelA) rectangle (8.40,\YmodelB);
    \draw[#1, opacity=0.40, line width=0.55pt, rounded corners=3pt]
      (1.05,\YmodelA) rectangle (8.40,\YmodelB);
  }

  \newcommand{\howbox}[1]{%
    \fill[#1, opacity=0.040, rounded corners=3pt] (0.55,\YhowA) rectangle (8.95,\YhowB);
    \draw[#1, opacity=0.75, line width=0.7pt, rounded corners=3pt]
      (0.55,\YhowA) rectangle (8.95,\YhowB);
  }

  \newcommand{\outbox}[1]{%
    \fill[#1, opacity=0.040, rounded corners=3pt] (0.55,\YoutA) rectangle (8.95,\YoutB);
    \draw[#1, opacity=0.75, line width=0.7pt, rounded corners=3pt]
      (0.55,\YoutA) rectangle (8.95,\YoutB);
  }

  \newcommand{\imageorfallback}[3]{%
    \IfFileExists{#1}{%
      \includegraphics[width=#2]{#1}%
    }{%
      \fbox{\parbox[c][2.8cm][c]{#2}{\centering\scriptsize #3}}%
    }%
  }

  % output grid with observed context and hidden interval
  \newcommand{\miniout}[1]{%
    \foreach \s in {1,...,9}{
      \draw[vgGrid,line width=0.28pt](0.2+\s*0.9,#1)--(0.2+\s*0.9,#1+3.3);
    }
    \fill[vgHidden,opacity=0.12](0.2+3.5*0.9,#1)rectangle(0.2+6.5*0.9,#1+3.3);
    \draw[vgTarget,line width=1.0pt](1.1,#1+1.35)--(2.0,#1+1.00)--(2.9,#1+1.20);
    \draw[vgTarget,line width=1.0pt](6.5,#1+2.50)--(7.4,#1+2.85)--(8.3,#1+2.50);
    \foreach \s/\y in {1/1.35,2/1.00,3/1.20,7/2.50,8/2.85,9/2.50}{
      \fill[vgTarget](0.2+\s*0.9,#1+\y)circle(0.10);
    }
  }

  % compact target-gap sketch used in "how" boxes
  \newcommand{\howgap}[1]{%
    \foreach \s in {1,...,9}{
      \draw[vgGrid,line width=0.25pt](0.2+\s*0.9,#1)--(0.2+\s*0.9,#1+3.3);
    }
    \fill[vgHidden,opacity=0.12](0.2+3.5*0.9,#1)rectangle(0.2+6.5*0.9,#1+3.3);
    \draw[vgTarget,line width=0.95pt](1.1,#1+1.35)--(2.0,#1+1.00)--(2.9,#1+1.20);
    \draw[vgTarget,line width=0.95pt](6.5,#1+2.50)--(7.4,#1+2.85)--(8.3,#1+2.50);
    \foreach \s/\y in {1/1.35,2/1.00,3/1.20,7/2.50,8/2.85,9/2.50}{
      \fill[vgTarget](0.2+\s*0.9,#1+\y)circle(0.09);
    }
  }

  % left row labels
  \node[rotate=90,font=\footnotesize\bfseries,text=vgFaint] at (-2.25,22.10) {METHOD};
  \node[rotate=90,font=\footnotesize\bfseries,text=vgFaint] at (-2.25,15.10) {MECHANISM};
  \node[rotate=90,font=\footnotesize\bfseries,text=vgFaint] at (-2.25,5.35) {OUTPUT};

  % ---------------------------------------------------------------------------
  % Panel 1: Simple baselines
  % ---------------------------------------------------------------------------
  \begin{scope}[shift={(0,0)}]
    \panelbg{Simple baselines}{mBase}
    \modelchip{mBase}
    \node[font=\tiny, text=vgInk, align=center, text width=3.2cm] at (4.72,21.65)
      {climatology\\persistence\\linear interpolation};

    \howbox{mBase}
    % three rule sketches with equations, not names
    \def\yshift{-0.9}
\foreach \yy/\eqn in {
  17.90/{{$\hat y_t = y_{t^-}$}},
  15.90/{{$\hat y_t=\tilde y_t^{\mathrm{lin}}$}},
  13.90/{{$\hat y_t = \mu_{\mathrm{season}}$}}
}{
  \fill[vgHidden,opacity=0.12](1.35,\yy-0.55+\yshift) rectangle (2.30,\yy+0.55+\yshift);
  \fill[vgTarget](0.95,\yy-0.25+\yshift) circle(0.09);
  \fill[vgTarget](2.80,\yy+0.30+\yshift) circle(0.09);
  \node[font=\scriptsize,text=vgInk,anchor=west] at (3.60,\yy+\yshift) {\eqn};
}
    % persistence
    \draw[mBase,line width=1.35pt] (0.95,16.75)--(2.80,16.75);
    % interpolation
    \draw[mBase,line width=1.35pt] (0.95,14.75)--(2.80,15.30);
    % climatology
    \draw[mBase,line width=1.35pt] (0.95,13)--(2.80,13);

    \draw[flow] (4.75,20.35)--(4.75,19.95);
    \draw[flow] (4.75,10.25)--(4.75,9.55);

    \outbox{mBase}
    \node[font=\scriptsize,text=vgInk,align=center] at (4.75,\Yeq)
      {$\hat y_t=\tilde y_t^{\mathrm{lin}}$};
    \miniout{\ob}
    % only interpolation shown in output
    \draw[mBase,line width=1.6pt](2.9,\ob+1.2)--(6.5,\ob+2.5);
    \foreach \s in {4,5,6}{
      \fill[mBase](0.2+\s*0.9,{\ob+1.2+(\s-3)*0.325})circle(0.10);
    }
  \end{scope}

  % ---------------------------------------------------------------------------
  % Panel 2: Target-only smoother
  % ---------------------------------------------------------------------------
  \begin{scope}[shift={(\DX,0)}]
    \panelbg{\scalebox{0.93}{Target-only smoother}}{mGP}
    \modelchip{mGP}
    \node[font=\scriptsize, text=vgInk, align=center, text width=3.0cm] at (4.72,21.65)
      {Gaussian process};

    \howbox{mGP}

\begin{scope}[yshift=-0.35cm]
\howgap{13.85}

% Prior mean m(t): faint dashed reference trend
\draw[vgFaint,line width=0.75pt,dash pattern=on 2pt off 1.5pt]
  plot[smooth] coordinates{
    (1.1,15.20)
    (2.9,15.10)
    (4.7,15.70)
    (6.5,16.35)
    (8.3,16.25)
  };
\node[font=\scriptsize,text=vgFaint,anchor=west] at (6.75,14.75) {$m(t)$};

% Posterior uncertainty induced by the covariance kernel inside/near the gap
\fill[mGP,opacity=0.13]
  plot[smooth] coordinates{
    (2.9,15.32)
    (3.6,15.55)
    (4.7,16.25)
    (5.8,16.55)
    (6.5,16.58)
  }
  -- plot[smooth] coordinates{
    (6.5,16.12)
    (5.8,15.55)
    (4.7,15.05)
    (3.6,14.92)
    (2.9,14.92)
  }
  -- cycle;

% Posterior mean reconstruction
\draw[mGP,line width=1.65pt]
  plot[smooth] coordinates{
    (2.9,15.05)
    (3.6,15.18)
    (4.7,15.75)
    (5.8,16.12)
    (6.5,16.35)
  };

% Small covariance cue
\draw[mGP,line width=0.55pt,opacity=0.85]
  (4.15,16.15) to[bend left=30] (5.25,16.15);
\draw[mGP,line width=0.55pt,opacity=0.65]
  (3.80,16.42) to[bend left=35] (5.60,16.42);
\node[font=\scriptsize,text=mGP] at (4.70,17.80) {$k(t,t')$};
\end{scope}

\draw[flow] (4.75,20.35)--(4.75,19.95);
\draw[flow] (4.75,10.25)--(4.75,9.55);

    \outbox{mGP}
    \node[font=\scriptsize,text=vgInk,align=center] at (4.75,\Yeq)
      {$\hat y(t)\sim\mathcal{GP}(m,k)$};
    \miniout{\ob}
    \fill[mGP,opacity=0.14]
      plot[smooth]coordinates{(2.9,\ob+1.25)(3.8,\ob+1.20)(4.7,\ob+2.75)(5.6,\ob+2.25)(6.5,\ob+2.55)}
      --plot[smooth]coordinates{(6.5,\ob+2.45)(5.6,\ob+1.55)(4.7,\ob+2.00)(3.8,\ob+0.60)(2.9,\ob+1.15)}
      --cycle;
    \draw[mGP,line width=1.6pt]
      plot[smooth]coordinates{(2.9,\ob+1.20)(3.8,\ob+0.90)(4.7,\ob+2.40)(5.6,\ob+1.90)(6.5,\ob+2.50)};
  \end{scope}

  % ---------------------------------------------------------------------------
  % Panel 3: External tabular
  % ---------------------------------------------------------------------------
  \begin{scope}[shift={({2*\DX},0)}]
    \panelbg{External tabular}{mExt}
    \modelchip{mExt}
    \node[font=\scriptsize, text=vgInk, align=center, text width=3.0cm] at (4.72,21.65)
      {HGB, ExtraTrees};

    \howbox{mExt}
    \node[inner sep=0pt] at (4.75,15.40)
      {\imageorfallback{figures/model_ladder_external_tabular.png}{3.35cm}{external tabular subfigure}};
    % symbol overlay
    \node[font=\scriptsize,text=mExt] at (6.40,12.6) {$\phi_t$};

    \draw[flow] (4.75,20.35)--(4.75,19.95);
    \draw[flow] (4.75,10.25)--(4.75,9.55);

    \outbox{mExt}
    \node[font=\scriptsize,text=vgInk,align=center] at (4.75,\Yeq)
      {$\hat y_t = f(\phi_t)$};
    \miniout{\ob}
    \foreach \s/\y in {4/1.35,5/1.90,6/1.55}{
      \fill[mExt](0.2+\s*0.9,\ob+\y)circle(0.15);
    }
  \end{scope}

  % ---------------------------------------------------------------------------
  % Panel 4: Gap-edge residual
  % ---------------------------------------------------------------------------
  \begin{scope}[shift={({3*\DX},0)}]
    \panelbg{Gap-edge residual}{mGap}
    \modelchip{mGap}
    \node[font=\scriptsize, text=vgInk, align=center, text width=3.0cm] at (4.72,21.65)
      {HGB, ExtraTrees};

    \howbox{mGap}
    \node[inner sep=0pt] at (4.75,15.40)
      {\imageorfallback{figures/model_ladder_gap_edge_residual.png}{3.30cm}{gap-edge residual subfigure}};
    % symbol overlays
    \node[font=\scriptsize,text=mExt] at (6.72,12.9) {$\phi_t$};
    \node[font=\scriptsize,text=mGap] at (1.7,12) {$\psi_G$};

    \draw[flow] (4.75,20.35)--(4.75,19.95);
    \draw[flow] (4.75,10.25)--(4.75,9.55);

    \outbox{mGap}
    \node[font=\scriptsize,text=vgInk,align=center] at (4.75,\Yeq)
      {$\hat y_t=\tilde y_t^{\mathrm{lin}} + g(\phi_t,\psi_G)$};
    \miniout{\ob}
    % interpolation baseline
    \draw[vgArrow,line width=1.1pt,dash pattern=on 1.7pt off 1.4pt](2.9,\ob+1.2)--(6.5,\ob+2.5);
    % residual bars only; no connecting curve
    \foreach \s/\by/\cy in {4/1.525/0.95,5/1.85/2.42,6/2.175/1.92}{
      \draw[mGap,line width=1.35pt](0.2+\s*0.9,\ob+\by)--(0.2+\s*0.9,\ob+\cy);
      \fill[mGap](0.2+\s*0.9,\ob+\cy)circle(0.11);
    }
  \end{scope}

  % ---------------------------------------------------------------------------
  % Panel 5: TS-ICL
  % ---------------------------------------------------------------------------
  \begin{scope}[shift={({4*\DX},0)}]
    \panelbg{TS-ICL}{mTS}
    \modelchip{mTS}
    \node[font=\scriptsize, text=vgInk, align=center, text width=3.0cm] at (4.72,21.65)
      {pretrained\\zero-shot};

    \howbox{mTS}
    \node[inner sep=0pt] at (4.75,14.9)
      {\imageorfallback{figures/model_ladder_tsicl.png}{2.75cm}{TS-ICL subfigure}};
    % symbol overlays
    \node[font=\scriptsize,text=tsgreen,anchor=west] at (1.3,19) {$ \mathbf y_{\mathrm{ctxt}}$};
    \node[font=\scriptsize,text=tsblue,anchor=west] at (3.1,16) {$\mathbf X_{\mathrm{covar}}$};

    \draw[flow] (4.75,20.35)--(4.75,19.95);
    \draw[flow] (4.75,10.25)--(4.75,9.55);

    \outbox{mTS}
    \node[font=\scriptsize,text=vgInk,align=center] at (4.75,\Yeq)
      {$ \hat y(t)= F(\mathbf y_{\mathrm{ctxt}}, \mathbf X_{\mathrm{covar}})$};
    \miniout{\ob}
    \fill[mTS,opacity=0.13]
      plot[smooth]coordinates{(2.9,\ob+1.25)(3.8,\ob+1.25)(4.7,\ob+2.55)(5.6,\ob+2.20)(6.5,\ob+2.55)}
      --plot[smooth]coordinates{(6.5,\ob+2.45)(5.6,\ob+1.70)(4.7,\ob+1.85)(3.8,\ob+0.85)(2.9,\ob+1.15)}
      --cycle;
    \draw[mTS,line width=1.6pt]
      plot[smooth]coordinates{(2.9,\ob+1.20)(3.8,\ob+1.05)(4.7,\ob+2.20)(5.6,\ob+1.95)(6.5,\ob+2.50)};
  \end{scope}

  % ---------------------------------------------------------------------------
  % Top arrow
  % ---------------------------------------------------------------------------
  \draw[-{Latex[length=3mm]},vgFaint,line width=1.35pt](0.2,25.45)--(51.35,25.45);
  \node[font=\small\itshape,text=vgInk] at (25.75,26.20)
    {increasing information and modelling structure};

\end{tikzpicture}}%
  }
  \caption{Comparator groups in the reconstruction benchmark.
  The five groups solve the same artificial-gap task but use different information structures:
  fixed temporal rules, target-only Gaussian-process smoothing, day-wise tabular prediction from
  external features, residual correction over interpolation using gap-edge context, and zero-shot sequence imputation.}
  \label{fig:model_ladder}
\end{figure}

\textbf{Simple baselines.}
The first group contains climatology, persistence, and linear interpolation. These methods have
no station-specific fitted model. Persistence fills the gap with the last available observation,
\[
    \hat y_t = y_{t^-},
\]
where $t^-$ is the closest observed time before the gap. Climatology predicts a typical seasonal
value, denoted here by $\mu_{\mathrm{season}}$. For an internal gap with observed endpoints
$t^-<t<t^+$, linear interpolation is
\[
    \tilde y_t^{\mathrm{lin}}
    =
    (1-\alpha_t)y_{t^-}+\alpha_t y_{t^+},
    \qquad
    \alpha_t=\frac{t-t^-}{t^+-t^-}.
\]
These baselines are deliberately simple. They define the level that more structured methods must
exceed, and linear interpolation is especially important because many observed gaps are short.
\vspace{0.25cm}

\textbf{Target-only smoother.}
The target-only smoother is a Gaussian-process model using visible target values around the gap
and no external covariates. In local notation, the target is treated as a realization of a latent
function
\[
    y(t) = f(t)+\varepsilon_t,
    \qquad
    f\sim\mathcal{GP}(m,k),
\]
with mean function $m(t)$ and Mat\'ern covariance kernel $k(t,t')$. If $\mathcal O_G$ denotes the
visible local context used by the smoother, the posterior mean over the hidden timestamps can be
written as
\[
    \hat{\mathbf y}_G
    =
    \mathbf m_G
    +
    K_{G,\mathcal O}
    \left(K_{\mathcal O,\mathcal O}+\sigma^2 I\right)^{-1}
    \left(\mathbf y_{\mathcal O}-\mathbf m_{\mathcal O}\right).
\]
This comparator tests how much of the gap can be recovered from local target continuity alone,
without satellite, meteorological, SST, or current predictors \citep{rasmussen2006gaussian}.
\vspace{0.25cm}

\textbf{External tabular models.}
The external tabular arm represents each hidden day by an engineered feature vector
\[
    \phi_t=\phi(t,\mathbf{x}),
\]
built from calendar, satellite, meteorological, SST, and current-related predictors. The notation
$\phi(t,\mathbf{x})$ includes the temporal and spatial summaries described in
Section~\ref{sec:predictors}; for example, lags, rolling summaries, regional statistics, and
nearby-grid-cell values. A supervised tabular regressor then predicts each hidden day separately:
\[
    \hat y_t = f(\phi_t).
\]
The benchmark uses two tree-ensemble learners for this arm, ExtraTrees and histogram gradient
boosting \citep{geurts2006extremely,friedman2001greedy,hastie2009elements}. This group tests
whether external covariates alone contain enough information to reconstruct the target.
\vspace{0.25cm}

\textbf{Gap-edge residual models.}
The gap-edge residual arm uses the same tabular learner family, but changes both the input and
the target of the supervised problem. It augments the external feature vector $\phi_t$ with
target-derived gap-edge summaries
\[
    \psi_G=\psi\!\left(\mathbf y_{\mathrm{edge}(G)}\right),
\]
computed from visible observations before and after the artificial gap, such as boundary values,
pre-/post-gap rolling means, and edge-to-edge slope. Rather than predicting $y_t$ directly, the
model predicts the residual over linear interpolation:
\[
    r_t = y_t-\tilde y_t^{\mathrm{lin}},
    \qquad
    \hat r_t = g(\phi_t,\psi_G),
    \qquad
    \hat y_t=\tilde y_t^{\mathrm{lin}}+\hat r_t .
\]
This formulation asks a narrower question than the external tabular arm. It does not try to
replace interpolation from scratch; it tests whether external features and gap-edge context can
detect systematic interpolation error.
\vspace{0.25cm}

\textbf{TS-ICL.}
TS-ICL is a pretrained time-series foundation model used here in zero-shot mode, with no
station-specific fitting or fine-tuning \citep{lenaour2026tsicl,tsicl_software}. For an
imputation gap, let $\mathcal T_{\mathrm{ctxt}}$ denote the visible target timestamps and
$\mathcal T_{\mathrm{tgt}}=G$ the hidden timestamps. The TS-ICL problem can be written as
conditional inference
\[
    p\!\left(\mathbf y_G
    \mid
    \mathbf y_{\mathrm{ctxt}},\mathbf X_{\mathrm{covar}}\right),
\]
where $\mathbf X_{\mathrm{covar}}$ denotes optional aligned or partially available covariate
channels. Architecturally, TS-ICL first encodes timestamp--value observations and covariates into
latent temporal representations, aggregates information across channels, builds
timestamp-specific context representations, and finally applies an in-context regressor based on
causal self-attention. In the notation of Figure~\ref{fig:model_ladder}, the point
reconstruction used for scoring is summarized as
\[
    \hat{\mathbf y}_G = F(\mathbf y_{\mathrm{ctxt}},\mathbf X_{\mathrm{covar}}).
\]
Unlike the tabular models, TS-ICL does not receive one fixed feature vector per hidden day. It
receives a masked target sequence and optional time-indexed covariate channels, and reconstructs
the hidden interval from sequence context. The covariate channels supplied to TS-ICL may still be
engineered physical summaries, but they enter as time-indexed channels rather than as a single
per-row tabular design matrix.

\subsection{Artificial gaps, real gaps, and validation scope}
\label{sec:artificial_gaps}

Methods are ranked only on artificial gaps: observed eligible target values are hidden,
reconstructed, and compared with the withheld truth. Real gaps have no recorded truth, so their
reconstructions can support later analyses that require continuous series, but they cannot
validate or rank methods.

For chlorophyll, the artificial-gap pool contains 681 gaps with lengths up to 60 days,
stratified by season and year. The 256-day chlorophyll scenario is kept outside the scored pool
and is used only to illustrate the station's longest real outage. For dissolved oxygen, the pool
contains 412 gaps: 406 primary gaps with lengths up to 30 days and six longer exploratory gaps.
Oxygen claims are therefore restricted to the primary $L \leq 30$ range. Longer gaps are harder
to validate because they require uninterrupted eligible observed runs; as $L$ increases, feasible
artificial gaps become fewer and long-gap estimates are less stable.

\begin{figure}[!htbp]
  \centering
  \resizebox{0.95\textwidth}{!}{% TikZ source for Figure 3 (R2 redesign): artificial-gap vs. real-gap validation scope.
% Included via \input from report/sections/03_methodology.tex, inside:
%   \resizebox{\textwidth}{!}{% TikZ source for Figure 3 (R2 redesign): artificial-gap vs. real-gap validation scope.
% Included via \input from report/sections/03_methodology.tex, inside:
%   \resizebox{\textwidth}{!}{% TikZ source for Figure 3 (R2 redesign): artificial-gap vs. real-gap validation scope.
% Included via \input from report/sections/03_methodology.tex, inside:
%   \resizebox{\textwidth}{!}{\input{figures/fig3_validation_schematic_tikz.tex}}
% Scope of THIS figure: only the distinction between (row 1) artificial-gap validation, where a
% withheld truth exists and reconstruction can be scored, and (row 2) real-gap reconstruction,
% where no truth exists and the result is not validation evidence. LOCO, TS-ICL context, tabular
% training rows and covariate availability are deliberately NOT shown here (separate figure).
% Storyboard grid: 2 rows (artificial / real) x 3 columns (recorded series -> reconstruction
% operation -> validation meaning). No packages needed beyond the report's existing TikZ setup
% (arrows.meta, decorations.pathreplacing); grey hatch is drawn with manual diagonals, not the
% patterns library.
\definecolor{vgObserved}{HTML}{2A7F8E}  % observed target (muted teal)
\definecolor{vgHidden}{HTML}{C0504D}    % artificial mask / withheld truth (muted red)
\definecolor{vgRecon}{HTML}{1B9E77}     % reconstruction (muted green)
\definecolor{vgMiss}{HTML}{9AA6B2}      % real missing segment (light grey hatch)
\definecolor{vgInk}{HTML}{2B2B2B}       % text / rules (dark grey)

\begin{tikzpicture}[x=0.5cm, y=0.5cm, font=\small,
    obs/.style={fill=vgObserved, draw=vgInk, line width=0.2pt},
    mask/.style={fill=vgHidden, draw=vgInk, line width=0.2pt},
    flow/.style={-{Latex[length=2.1mm]}, vgInk, line width=0.7pt},
    recon/.style={vgRecon, line width=1.3pt, dash pattern=on 3pt off 2pt},
    rdot/.style={circle, fill=vgRecon, draw=none, inner sep=1.0pt}]

  % --- geometry ---------------------------------------------------------------
  % columns: X1 recorded series, X2 reconstruction operation, X3 validation meaning
  % rows:    Y1 artificial gap (top), Y2 real gap (bottom); cell w=0.8, h=0.7, step 1.0
  \def\XA{0}   \def\XB{13}  \def\XC{25.4}
  \def\YA{3.4} \def\YB{0.0}
  % centres (cell centre y = row + 0.35; strip centre x = origin + 4.9)

  % --- column headers ---------------------------------------------------------
  \node[font=\bfseries, text=vgInk] at (4.9,5.75)  {Recorded series};
  \node[font=\bfseries, text=vgInk] at (17.9,5.75) {Reconstruction operation};
  \node[font=\bfseries, text=vgInk, anchor=west] at (24.6,5.75) {Validation meaning};

  % --- row labels -------------------------------------------------------------
  \node[font=\bfseries, text=vgInk, align=center] at (-2.1,3.75) {Artificial\\gap};
  \node[font=\bfseries, text=vgInk, align=center] at (-2.1,0.35) {Real\\gap};

  % faint separator between the two storyboard rows
  \draw[vgMiss, line width=0.3pt] (-3.4,1.95) -- (33.6,1.95);

  % ===========================================================================
  % ROW 1 -- ARTIFICIAL GAP
  % ===========================================================================
  % Col 1: observed series, with the segment that will be masked marked as known values
  \foreach \i in {0,...,9}{ \fill[obs] (\XA+\i,\YA) rectangle (\XA+\i+0.8,\YA+0.7); }
  \draw[vgHidden, line width=0.7pt, dash pattern=on 2pt off 1.5pt]
      (\XA+3.9,\YA-0.12) rectangle (\XA+6.9,\YA+0.82);
  \draw[decorate,decoration={brace,amplitude=3pt},line width=0.4pt,vgInk]
      (\XA+3.9,\YA+0.95) -- (\XA+6.9,\YA+0.95)
      node[midway,above=1pt,font=\footnotesize,text=vgInk] {known values};
  \node[font=\small, text=vgInk] at (4.9,\YA-0.75) {observed values};

  % Col 2: hide the known segment (red), then reconstruct (green)
  \foreach \i in {0,1,2,3,7,8,9}{ \fill[obs]  (\XB+\i,\YA) rectangle (\XB+\i+0.8,\YA+0.7); }
  \foreach \i in {4,5,6}{        \fill[mask] (\XB+\i,\YA) rectangle (\XB+\i+0.8,\YA+0.7); }
  \draw[recon] plot[smooth] coordinates
      {(\XB+3.9,\YA+0.35) (\XB+4.4,\YA+0.58) (\XB+5.4,\YA+0.18)
       (\XB+6.4,\YA+0.55) (\XB+6.9,\YA+0.33)};
  \foreach \x/\y in {4.4/0.58, 5.4/0.18, 6.4/0.55}{ \node[rdot] at (\XB+\x,\YA+\y) {}; }
  \node[font=\footnotesize, text=vgHidden] at (\XB+5.0,\YA+1.15) {hide known segment};
  \node[font=\footnotesize, text=vgRecon]  at (\XB+5.0,\YA-0.55) {reconstruct};

  % Col 3: truth exists -> score error
  \node[anchor=west, align=left, text=vgInk] at (\XC-0.25,\YA+0.35)
      {compare reconstruction with\\ hidden observed values\\
       $\Rightarrow$ compute MAE};

  % ===========================================================================
  % ROW 2 -- REAL GAP
  % ===========================================================================
  % Col 1: observed series with a genuine missing segment (grey hatch, no data)
  \foreach \i in {0,1,2,3,7,8,9}{ \fill[obs] (\XA+\i,\YB) rectangle (\XA+\i+0.8,\YB+0.7); }
  \foreach \i in {4,5,6}{
      \draw[draw=vgMiss,line width=0.4pt] (\XA+\i,\YB) rectangle (\XA+\i+0.8,\YB+0.7);
      \draw[vgMiss,line width=0.4pt] (\XA+\i,\YB)      -- (\XA+\i+0.8,\YB+0.7);
      \draw[vgMiss,line width=0.4pt] (\XA+\i+0.4,\YB)  -- (\XA+\i+0.8,\YB+0.35);
      \draw[vgMiss,line width=0.4pt] (\XA+\i,\YB+0.35) -- (\XA+\i+0.4,\YB+0.7);
  }
  \node[font=\small, text=vgInk] at (4.9,\YB-0.75) {real missing segment};

  % Col 2: reconstruction fills the blank gap (no underlying truth to expose)
  \foreach \i in {0,1,2,3,7,8,9}{ \fill[obs] (\XB+\i,\YB) rectangle (\XB+\i+0.8,\YB+0.7); }
  \foreach \i in {4,5,6}{
      \draw[draw=vgMiss,line width=0.4pt] (\XB+\i,\YB) rectangle (\XB+\i+0.8,\YB+0.7);
      \draw[vgMiss,line width=0.4pt] (\XB+\i,\YB)      -- (\XB+\i+0.8,\YB+0.7);
      \draw[vgMiss,line width=0.4pt] (\XB+\i+0.4,\YB)  -- (\XB+\i+0.8,\YB+0.35);
      \draw[vgMiss,line width=0.4pt] (\XB+\i,\YB+0.35) -- (\XB+\i+0.4,\YB+0.7);
  }
  \draw[recon] plot[smooth] coordinates
      {(\XB+3.9,\YB+0.35) (\XB+4.4,\YB+0.58) (\XB+5.4,\YB+0.18)
       (\XB+6.4,\YB+0.55) (\XB+6.9,\YB+0.33)};
  \foreach \x/\y in {4.4/0.58, 5.4/0.18, 6.4/0.55}{ \node[rdot] at (\XB+\x,\YB+\y) {}; }
  \node[font=\footnotesize, text=vgRecon] at (\XB+5.0,\YB-0.55) {reconstruct};

  % Col 3: no truth -> no score
  \node[anchor=west, align=left, text=vgInk] at (\XC-0.25,\YB+0.35)
      {nothing to compare against\\
       $\Rightarrow$ not validation evidence};
  % ===========================================================================
  % INTER-COLUMN FLOW ARROWS (both rows)
  % ===========================================================================
  \draw[flow] (10.2,\YA+0.35) -- (12.7,\YA+0.35);
  \draw[flow] (10.2,\YB+0.35) -- (12.7,\YB+0.35);
  \draw[flow] (23.1,\YA+0.35) -- (24.7,\YA+0.35);
  \draw[flow] (23.1,\YB+0.35) -- (24.7,\YB+0.35);

\end{tikzpicture}
}
% Scope of THIS figure: only the distinction between (row 1) artificial-gap validation, where a
% withheld truth exists and reconstruction can be scored, and (row 2) real-gap reconstruction,
% where no truth exists and the result is not validation evidence. LOCO, TS-ICL context, tabular
% training rows and covariate availability are deliberately NOT shown here (separate figure).
% Storyboard grid: 2 rows (artificial / real) x 3 columns (recorded series -> reconstruction
% operation -> validation meaning). No packages needed beyond the report's existing TikZ setup
% (arrows.meta, decorations.pathreplacing); grey hatch is drawn with manual diagonals, not the
% patterns library.
\definecolor{vgObserved}{HTML}{2A7F8E}  % observed target (muted teal)
\definecolor{vgHidden}{HTML}{C0504D}    % artificial mask / withheld truth (muted red)
\definecolor{vgRecon}{HTML}{1B9E77}     % reconstruction (muted green)
\definecolor{vgMiss}{HTML}{9AA6B2}      % real missing segment (light grey hatch)
\definecolor{vgInk}{HTML}{2B2B2B}       % text / rules (dark grey)

\begin{tikzpicture}[x=0.5cm, y=0.5cm, font=\small,
    obs/.style={fill=vgObserved, draw=vgInk, line width=0.2pt},
    mask/.style={fill=vgHidden, draw=vgInk, line width=0.2pt},
    flow/.style={-{Latex[length=2.1mm]}, vgInk, line width=0.7pt},
    recon/.style={vgRecon, line width=1.3pt, dash pattern=on 3pt off 2pt},
    rdot/.style={circle, fill=vgRecon, draw=none, inner sep=1.0pt}]

  % --- geometry ---------------------------------------------------------------
  % columns: X1 recorded series, X2 reconstruction operation, X3 validation meaning
  % rows:    Y1 artificial gap (top), Y2 real gap (bottom); cell w=0.8, h=0.7, step 1.0
  \def\XA{0}   \def\XB{13}  \def\XC{25.4}
  \def\YA{3.4} \def\YB{0.0}
  % centres (cell centre y = row + 0.35; strip centre x = origin + 4.9)

  % --- column headers ---------------------------------------------------------
  \node[font=\bfseries, text=vgInk] at (4.9,5.75)  {Recorded series};
  \node[font=\bfseries, text=vgInk] at (17.9,5.75) {Reconstruction operation};
  \node[font=\bfseries, text=vgInk, anchor=west] at (24.6,5.75) {Validation meaning};

  % --- row labels -------------------------------------------------------------
  \node[font=\bfseries, text=vgInk, align=center] at (-2.1,3.75) {Artificial\\gap};
  \node[font=\bfseries, text=vgInk, align=center] at (-2.1,0.35) {Real\\gap};

  % faint separator between the two storyboard rows
  \draw[vgMiss, line width=0.3pt] (-3.4,1.95) -- (33.6,1.95);

  % ===========================================================================
  % ROW 1 -- ARTIFICIAL GAP
  % ===========================================================================
  % Col 1: observed series, with the segment that will be masked marked as known values
  \foreach \i in {0,...,9}{ \fill[obs] (\XA+\i,\YA) rectangle (\XA+\i+0.8,\YA+0.7); }
  \draw[vgHidden, line width=0.7pt, dash pattern=on 2pt off 1.5pt]
      (\XA+3.9,\YA-0.12) rectangle (\XA+6.9,\YA+0.82);
  \draw[decorate,decoration={brace,amplitude=3pt},line width=0.4pt,vgInk]
      (\XA+3.9,\YA+0.95) -- (\XA+6.9,\YA+0.95)
      node[midway,above=1pt,font=\footnotesize,text=vgInk] {known values};
  \node[font=\small, text=vgInk] at (4.9,\YA-0.75) {observed values};

  % Col 2: hide the known segment (red), then reconstruct (green)
  \foreach \i in {0,1,2,3,7,8,9}{ \fill[obs]  (\XB+\i,\YA) rectangle (\XB+\i+0.8,\YA+0.7); }
  \foreach \i in {4,5,6}{        \fill[mask] (\XB+\i,\YA) rectangle (\XB+\i+0.8,\YA+0.7); }
  \draw[recon] plot[smooth] coordinates
      {(\XB+3.9,\YA+0.35) (\XB+4.4,\YA+0.58) (\XB+5.4,\YA+0.18)
       (\XB+6.4,\YA+0.55) (\XB+6.9,\YA+0.33)};
  \foreach \x/\y in {4.4/0.58, 5.4/0.18, 6.4/0.55}{ \node[rdot] at (\XB+\x,\YA+\y) {}; }
  \node[font=\footnotesize, text=vgHidden] at (\XB+5.0,\YA+1.15) {hide known segment};
  \node[font=\footnotesize, text=vgRecon]  at (\XB+5.0,\YA-0.55) {reconstruct};

  % Col 3: truth exists -> score error
  \node[anchor=west, align=left, text=vgInk] at (\XC-0.25,\YA+0.35)
      {compare reconstruction with\\ hidden observed values\\
       $\Rightarrow$ compute MAE};

  % ===========================================================================
  % ROW 2 -- REAL GAP
  % ===========================================================================
  % Col 1: observed series with a genuine missing segment (grey hatch, no data)
  \foreach \i in {0,1,2,3,7,8,9}{ \fill[obs] (\XA+\i,\YB) rectangle (\XA+\i+0.8,\YB+0.7); }
  \foreach \i in {4,5,6}{
      \draw[draw=vgMiss,line width=0.4pt] (\XA+\i,\YB) rectangle (\XA+\i+0.8,\YB+0.7);
      \draw[vgMiss,line width=0.4pt] (\XA+\i,\YB)      -- (\XA+\i+0.8,\YB+0.7);
      \draw[vgMiss,line width=0.4pt] (\XA+\i+0.4,\YB)  -- (\XA+\i+0.8,\YB+0.35);
      \draw[vgMiss,line width=0.4pt] (\XA+\i,\YB+0.35) -- (\XA+\i+0.4,\YB+0.7);
  }
  \node[font=\small, text=vgInk] at (4.9,\YB-0.75) {real missing segment};

  % Col 2: reconstruction fills the blank gap (no underlying truth to expose)
  \foreach \i in {0,1,2,3,7,8,9}{ \fill[obs] (\XB+\i,\YB) rectangle (\XB+\i+0.8,\YB+0.7); }
  \foreach \i in {4,5,6}{
      \draw[draw=vgMiss,line width=0.4pt] (\XB+\i,\YB) rectangle (\XB+\i+0.8,\YB+0.7);
      \draw[vgMiss,line width=0.4pt] (\XB+\i,\YB)      -- (\XB+\i+0.8,\YB+0.7);
      \draw[vgMiss,line width=0.4pt] (\XB+\i+0.4,\YB)  -- (\XB+\i+0.8,\YB+0.35);
      \draw[vgMiss,line width=0.4pt] (\XB+\i,\YB+0.35) -- (\XB+\i+0.4,\YB+0.7);
  }
  \draw[recon] plot[smooth] coordinates
      {(\XB+3.9,\YB+0.35) (\XB+4.4,\YB+0.58) (\XB+5.4,\YB+0.18)
       (\XB+6.4,\YB+0.55) (\XB+6.9,\YB+0.33)};
  \foreach \x/\y in {4.4/0.58, 5.4/0.18, 6.4/0.55}{ \node[rdot] at (\XB+\x,\YB+\y) {}; }
  \node[font=\footnotesize, text=vgRecon] at (\XB+5.0,\YB-0.55) {reconstruct};

  % Col 3: no truth -> no score
  \node[anchor=west, align=left, text=vgInk] at (\XC-0.25,\YB+0.35)
      {nothing to compare against\\
       $\Rightarrow$ not validation evidence};
  % ===========================================================================
  % INTER-COLUMN FLOW ARROWS (both rows)
  % ===========================================================================
  \draw[flow] (10.2,\YA+0.35) -- (12.7,\YA+0.35);
  \draw[flow] (10.2,\YB+0.35) -- (12.7,\YB+0.35);
  \draw[flow] (23.1,\YA+0.35) -- (24.7,\YA+0.35);
  \draw[flow] (23.1,\YB+0.35) -- (24.7,\YB+0.35);

\end{tikzpicture}
}
% Scope of THIS figure: only the distinction between (row 1) artificial-gap validation, where a
% withheld truth exists and reconstruction can be scored, and (row 2) real-gap reconstruction,
% where no truth exists and the result is not validation evidence. LOCO, TS-ICL context, tabular
% training rows and covariate availability are deliberately NOT shown here (separate figure).
% Storyboard grid: 2 rows (artificial / real) x 3 columns (recorded series -> reconstruction
% operation -> validation meaning). No packages needed beyond the report's existing TikZ setup
% (arrows.meta, decorations.pathreplacing); grey hatch is drawn with manual diagonals, not the
% patterns library.
\definecolor{vgObserved}{HTML}{2A7F8E}  % observed target (muted teal)
\definecolor{vgHidden}{HTML}{C0504D}    % artificial mask / withheld truth (muted red)
\definecolor{vgRecon}{HTML}{1B9E77}     % reconstruction (muted green)
\definecolor{vgMiss}{HTML}{9AA6B2}      % real missing segment (light grey hatch)
\definecolor{vgInk}{HTML}{2B2B2B}       % text / rules (dark grey)

\begin{tikzpicture}[x=0.5cm, y=0.5cm, font=\small,
    obs/.style={fill=vgObserved, draw=vgInk, line width=0.2pt},
    mask/.style={fill=vgHidden, draw=vgInk, line width=0.2pt},
    flow/.style={-{Latex[length=2.1mm]}, vgInk, line width=0.7pt},
    recon/.style={vgRecon, line width=1.3pt, dash pattern=on 3pt off 2pt},
    rdot/.style={circle, fill=vgRecon, draw=none, inner sep=1.0pt}]

  % --- geometry ---------------------------------------------------------------
  % columns: X1 recorded series, X2 reconstruction operation, X3 validation meaning
  % rows:    Y1 artificial gap (top), Y2 real gap (bottom); cell w=0.8, h=0.7, step 1.0
  \def\XA{0}   \def\XB{13}  \def\XC{25.4}
  \def\YA{3.4} \def\YB{0.0}
  % centres (cell centre y = row + 0.35; strip centre x = origin + 4.9)

  % --- column headers ---------------------------------------------------------
  \node[font=\bfseries, text=vgInk] at (4.9,5.75)  {Recorded series};
  \node[font=\bfseries, text=vgInk] at (17.9,5.75) {Reconstruction operation};
  \node[font=\bfseries, text=vgInk, anchor=west] at (24.6,5.75) {Validation meaning};

  % --- row labels -------------------------------------------------------------
  \node[font=\bfseries, text=vgInk, align=center] at (-2.1,3.75) {Artificial\\gap};
  \node[font=\bfseries, text=vgInk, align=center] at (-2.1,0.35) {Real\\gap};

  % faint separator between the two storyboard rows
  \draw[vgMiss, line width=0.3pt] (-3.4,1.95) -- (33.6,1.95);

  % ===========================================================================
  % ROW 1 -- ARTIFICIAL GAP
  % ===========================================================================
  % Col 1: observed series, with the segment that will be masked marked as known values
  \foreach \i in {0,...,9}{ \fill[obs] (\XA+\i,\YA) rectangle (\XA+\i+0.8,\YA+0.7); }
  \draw[vgHidden, line width=0.7pt, dash pattern=on 2pt off 1.5pt]
      (\XA+3.9,\YA-0.12) rectangle (\XA+6.9,\YA+0.82);
  \draw[decorate,decoration={brace,amplitude=3pt},line width=0.4pt,vgInk]
      (\XA+3.9,\YA+0.95) -- (\XA+6.9,\YA+0.95)
      node[midway,above=1pt,font=\footnotesize,text=vgInk] {known values};
  \node[font=\small, text=vgInk] at (4.9,\YA-0.75) {observed values};

  % Col 2: hide the known segment (red), then reconstruct (green)
  \foreach \i in {0,1,2,3,7,8,9}{ \fill[obs]  (\XB+\i,\YA) rectangle (\XB+\i+0.8,\YA+0.7); }
  \foreach \i in {4,5,6}{        \fill[mask] (\XB+\i,\YA) rectangle (\XB+\i+0.8,\YA+0.7); }
  \draw[recon] plot[smooth] coordinates
      {(\XB+3.9,\YA+0.35) (\XB+4.4,\YA+0.58) (\XB+5.4,\YA+0.18)
       (\XB+6.4,\YA+0.55) (\XB+6.9,\YA+0.33)};
  \foreach \x/\y in {4.4/0.58, 5.4/0.18, 6.4/0.55}{ \node[rdot] at (\XB+\x,\YA+\y) {}; }
  \node[font=\footnotesize, text=vgHidden] at (\XB+5.0,\YA+1.15) {hide known segment};
  \node[font=\footnotesize, text=vgRecon]  at (\XB+5.0,\YA-0.55) {reconstruct};

  % Col 3: truth exists -> score error
  \node[anchor=west, align=left, text=vgInk] at (\XC-0.25,\YA+0.35)
      {compare reconstruction with\\ hidden observed values\\
       $\Rightarrow$ compute MAE};

  % ===========================================================================
  % ROW 2 -- REAL GAP
  % ===========================================================================
  % Col 1: observed series with a genuine missing segment (grey hatch, no data)
  \foreach \i in {0,1,2,3,7,8,9}{ \fill[obs] (\XA+\i,\YB) rectangle (\XA+\i+0.8,\YB+0.7); }
  \foreach \i in {4,5,6}{
      \draw[draw=vgMiss,line width=0.4pt] (\XA+\i,\YB) rectangle (\XA+\i+0.8,\YB+0.7);
      \draw[vgMiss,line width=0.4pt] (\XA+\i,\YB)      -- (\XA+\i+0.8,\YB+0.7);
      \draw[vgMiss,line width=0.4pt] (\XA+\i+0.4,\YB)  -- (\XA+\i+0.8,\YB+0.35);
      \draw[vgMiss,line width=0.4pt] (\XA+\i,\YB+0.35) -- (\XA+\i+0.4,\YB+0.7);
  }
  \node[font=\small, text=vgInk] at (4.9,\YB-0.75) {real missing segment};

  % Col 2: reconstruction fills the blank gap (no underlying truth to expose)
  \foreach \i in {0,1,2,3,7,8,9}{ \fill[obs] (\XB+\i,\YB) rectangle (\XB+\i+0.8,\YB+0.7); }
  \foreach \i in {4,5,6}{
      \draw[draw=vgMiss,line width=0.4pt] (\XB+\i,\YB) rectangle (\XB+\i+0.8,\YB+0.7);
      \draw[vgMiss,line width=0.4pt] (\XB+\i,\YB)      -- (\XB+\i+0.8,\YB+0.7);
      \draw[vgMiss,line width=0.4pt] (\XB+\i+0.4,\YB)  -- (\XB+\i+0.8,\YB+0.35);
      \draw[vgMiss,line width=0.4pt] (\XB+\i,\YB+0.35) -- (\XB+\i+0.4,\YB+0.7);
  }
  \draw[recon] plot[smooth] coordinates
      {(\XB+3.9,\YB+0.35) (\XB+4.4,\YB+0.58) (\XB+5.4,\YB+0.18)
       (\XB+6.4,\YB+0.55) (\XB+6.9,\YB+0.33)};
  \foreach \x/\y in {4.4/0.58, 5.4/0.18, 6.4/0.55}{ \node[rdot] at (\XB+\x,\YB+\y) {}; }
  \node[font=\footnotesize, text=vgRecon] at (\XB+5.0,\YB-0.55) {reconstruct};

  % Col 3: no truth -> no score
  \node[anchor=west, align=left, text=vgInk] at (\XC-0.25,\YB+0.35)
      {nothing to compare against\\
       $\Rightarrow$ not validation evidence};
  % ===========================================================================
  % INTER-COLUMN FLOW ARROWS (both rows)
  % ===========================================================================
  \draw[flow] (10.2,\YA+0.35) -- (12.7,\YA+0.35);
  \draw[flow] (10.2,\YB+0.35) -- (12.7,\YB+0.35);
  \draw[flow] (23.1,\YA+0.35) -- (24.7,\YA+0.35);
  \draw[flow] (23.1,\YB+0.35) -- (24.7,\YB+0.35);

\end{tikzpicture}
}
  \caption{Artificial-gap validation versus real-gap reconstruction.}
  \label{fig:validation_schematic}
\end{figure}

\subsection{Leakage control for target-derived features}
\label{sec:loco}

For tabular models, removing only the hidden target days is not always enough. Some input columns
are computed from the target series itself, such as lagged target values, rolling target means,
pre-/post-gap boundary statistics, and interpolation-residual features. A row outside the hidden
interval can therefore still leak information if its target-derived feature window reaches into
the artificial gap. Rows whose target-derived windows overlap the hidden interval are excluded from training; rows whose windows stay fully outside it are kept (Figure~\ref{fig:loco_dependency_window}).

\begin{figure}[!htbp]
  \centering
  \resizebox{0.6\textwidth}{!}{% TikZ source: leakage control -- fixed target-derived windows vs. hidden interval.
% Scope: a tabular training row is excluded when its fixed target-derived
% feature window overlaps the hidden target interval.
%
% Included via:
%   \resizebox{0.78\textwidth}{!}{% TikZ source: leakage control -- fixed target-derived windows vs. hidden interval.
% Scope: a tabular training row is excluded when its fixed target-derived
% feature window overlaps the hidden target interval.
%
% Included via:
%   \resizebox{0.78\textwidth}{!}{% TikZ source: leakage control -- fixed target-derived windows vs. hidden interval.
% Scope: a tabular training row is excluded when its fixed target-derived
% feature window overlaps the hidden target interval.
%
% Included via:
%   \resizebox{0.78\textwidth}{!}{\input{figures/fig_loco_dependency_window_tikz.tex}}

\definecolor{vgObserved}{HTML}{2A7F8E}  % observed target (muted teal)
\definecolor{vgHidden}{HTML}{C0504D}    % hidden target interval (muted red)
\definecolor{vgKept}{HTML}{3B6FA0}      % kept window (muted blue)
\definecolor{vgExcl}{HTML}{D98A2B}      % excluded window (muted orange)
\definecolor{vgInk}{HTML}{2B2B2B}       % text (dark grey)

\begin{tikzpicture}[x=0.42cm, y=0.42cm, font=\small,
    obs/.style={fill=vgObserved, draw=vgInk, line width=0.2pt},
    hid/.style={fill=vgHidden, draw=vgInk, line width=0.2pt}]

  % Hidden interval occupies timeline cells 6, 7, 8.
  \def\HidLo{5}
  \def\HidHi{9}
  \def\YT{4.0}

  % ---------------------------------------------------------------------------
  % Top target timeline
  % ---------------------------------------------------------------------------
  \foreach \i in {1,2,3,4,5,9,10,11,12,13}{
    \fill[obs] (\i,\YT) rectangle (\i+0.8,\YT+0.62);
  }

  \foreach \i in {6,7,8}{
    \fill[hid] (\i,\YT) rectangle (\i+0.8,\YT+0.62);
  }

  \draw[decorate,decoration={brace,amplitude=3pt},line width=0.4pt,vgHidden]
    (6,\YT+0.78) -- (8.8,\YT+0.78)
    node[midway,above=2pt,font=\scriptsize,text=vgHidden]
    {hidden target interval};

  % ---------------------------------------------------------------------------
  % Equal-length sliding window macro
  % ---------------------------------------------------------------------------
  \newcommand{\drawwin}[3]{%
    \foreach \k in {0,1,2,3,4}{%
      \fill[#3, fill opacity=0.22, draw=#3, line width=0.45pt]
        (#1+\k,#2) rectangle (#1+\k+0.8,#2+0.62);%
      \pgfmathtruncatemacro{\cell}{#1+\k}%
      \ifnum\cell>\HidLo
        \ifnum\cell<\HidHi
          \fill[#3, fill opacity=0.40]
            (#1+\k,#2) rectangle (#1+\k+0.8,#2+0.62);
          \draw[vgHidden, line width=0.35pt]
            (#1+\k,#2) -- (#1+\k+0.8,#2+0.62);
          \draw[vgHidden, line width=0.35pt]
            (#1+\k+0.4,#2) -- (#1+\k+0.8,#2+0.31);
          \draw[vgHidden, line width=0.35pt]
            (#1+\k,#2+0.31) -- (#1+\k+0.4,#2+0.62);
        \fi
      \fi
    }%
  }

  % ---------------------------------------------------------------------------
  % Four same-length target-derived windows
  % ---------------------------------------------------------------------------
  \drawwin{1}{2.90}{vgKept}   % clear of hidden interval -> kept
  \drawwin{4}{1.95}{vgExcl}   % overlaps hidden interval -> excluded
  \drawwin{6}{1.00}{vgExcl}   % overlaps hidden interval -> excluded
  \drawwin{9}{0.05}{vgKept}   % clear of hidden interval -> kept

  % ---------------------------------------------------------------------------
  % Verdict labels
  % ---------------------------------------------------------------------------
  \node[anchor=west, font=\scriptsize\bfseries, text=vgKept]
    at (13.9,3.20) {kept};
  \node[anchor=west, font=\scriptsize\bfseries, text=vgExcl]
    at (13.9,2.25) {excluded};
  \node[anchor=west, font=\scriptsize\bfseries, text=vgExcl]
    at (13.9,1.30) {excluded};
  \node[anchor=west, font=\scriptsize\bfseries, text=vgKept]
    at (13.9,0.35) {kept};

\end{tikzpicture}}

\definecolor{vgObserved}{HTML}{2A7F8E}  % observed target (muted teal)
\definecolor{vgHidden}{HTML}{C0504D}    % hidden target interval (muted red)
\definecolor{vgKept}{HTML}{3B6FA0}      % kept window (muted blue)
\definecolor{vgExcl}{HTML}{D98A2B}      % excluded window (muted orange)
\definecolor{vgInk}{HTML}{2B2B2B}       % text (dark grey)

\begin{tikzpicture}[x=0.42cm, y=0.42cm, font=\small,
    obs/.style={fill=vgObserved, draw=vgInk, line width=0.2pt},
    hid/.style={fill=vgHidden, draw=vgInk, line width=0.2pt}]

  % Hidden interval occupies timeline cells 6, 7, 8.
  \def\HidLo{5}
  \def\HidHi{9}
  \def\YT{4.0}

  % ---------------------------------------------------------------------------
  % Top target timeline
  % ---------------------------------------------------------------------------
  \foreach \i in {1,2,3,4,5,9,10,11,12,13}{
    \fill[obs] (\i,\YT) rectangle (\i+0.8,\YT+0.62);
  }

  \foreach \i in {6,7,8}{
    \fill[hid] (\i,\YT) rectangle (\i+0.8,\YT+0.62);
  }

  \draw[decorate,decoration={brace,amplitude=3pt},line width=0.4pt,vgHidden]
    (6,\YT+0.78) -- (8.8,\YT+0.78)
    node[midway,above=2pt,font=\scriptsize,text=vgHidden]
    {hidden target interval};

  % ---------------------------------------------------------------------------
  % Equal-length sliding window macro
  % ---------------------------------------------------------------------------
  \newcommand{\drawwin}[3]{%
    \foreach \k in {0,1,2,3,4}{%
      \fill[#3, fill opacity=0.22, draw=#3, line width=0.45pt]
        (#1+\k,#2) rectangle (#1+\k+0.8,#2+0.62);%
      \pgfmathtruncatemacro{\cell}{#1+\k}%
      \ifnum\cell>\HidLo
        \ifnum\cell<\HidHi
          \fill[#3, fill opacity=0.40]
            (#1+\k,#2) rectangle (#1+\k+0.8,#2+0.62);
          \draw[vgHidden, line width=0.35pt]
            (#1+\k,#2) -- (#1+\k+0.8,#2+0.62);
          \draw[vgHidden, line width=0.35pt]
            (#1+\k+0.4,#2) -- (#1+\k+0.8,#2+0.31);
          \draw[vgHidden, line width=0.35pt]
            (#1+\k,#2+0.31) -- (#1+\k+0.4,#2+0.62);
        \fi
      \fi
    }%
  }

  % ---------------------------------------------------------------------------
  % Four same-length target-derived windows
  % ---------------------------------------------------------------------------
  \drawwin{1}{2.90}{vgKept}   % clear of hidden interval -> kept
  \drawwin{4}{1.95}{vgExcl}   % overlaps hidden interval -> excluded
  \drawwin{6}{1.00}{vgExcl}   % overlaps hidden interval -> excluded
  \drawwin{9}{0.05}{vgKept}   % clear of hidden interval -> kept

  % ---------------------------------------------------------------------------
  % Verdict labels
  % ---------------------------------------------------------------------------
  \node[anchor=west, font=\scriptsize\bfseries, text=vgKept]
    at (13.9,3.20) {kept};
  \node[anchor=west, font=\scriptsize\bfseries, text=vgExcl]
    at (13.9,2.25) {excluded};
  \node[anchor=west, font=\scriptsize\bfseries, text=vgExcl]
    at (13.9,1.30) {excluded};
  \node[anchor=west, font=\scriptsize\bfseries, text=vgKept]
    at (13.9,0.35) {kept};

\end{tikzpicture}}

\definecolor{vgObserved}{HTML}{2A7F8E}  % observed target (muted teal)
\definecolor{vgHidden}{HTML}{C0504D}    % hidden target interval (muted red)
\definecolor{vgKept}{HTML}{3B6FA0}      % kept window (muted blue)
\definecolor{vgExcl}{HTML}{D98A2B}      % excluded window (muted orange)
\definecolor{vgInk}{HTML}{2B2B2B}       % text (dark grey)

\begin{tikzpicture}[x=0.42cm, y=0.42cm, font=\small,
    obs/.style={fill=vgObserved, draw=vgInk, line width=0.2pt},
    hid/.style={fill=vgHidden, draw=vgInk, line width=0.2pt}]

  % Hidden interval occupies timeline cells 6, 7, 8.
  \def\HidLo{5}
  \def\HidHi{9}
  \def\YT{4.0}

  % ---------------------------------------------------------------------------
  % Top target timeline
  % ---------------------------------------------------------------------------
  \foreach \i in {1,2,3,4,5,9,10,11,12,13}{
    \fill[obs] (\i,\YT) rectangle (\i+0.8,\YT+0.62);
  }

  \foreach \i in {6,7,8}{
    \fill[hid] (\i,\YT) rectangle (\i+0.8,\YT+0.62);
  }

  \draw[decorate,decoration={brace,amplitude=3pt},line width=0.4pt,vgHidden]
    (6,\YT+0.78) -- (8.8,\YT+0.78)
    node[midway,above=2pt,font=\scriptsize,text=vgHidden]
    {hidden target interval};

  % ---------------------------------------------------------------------------
  % Equal-length sliding window macro
  % ---------------------------------------------------------------------------
  \newcommand{\drawwin}[3]{%
    \foreach \k in {0,1,2,3,4}{%
      \fill[#3, fill opacity=0.22, draw=#3, line width=0.45pt]
        (#1+\k,#2) rectangle (#1+\k+0.8,#2+0.62);%
      \pgfmathtruncatemacro{\cell}{#1+\k}%
      \ifnum\cell>\HidLo
        \ifnum\cell<\HidHi
          \fill[#3, fill opacity=0.40]
            (#1+\k,#2) rectangle (#1+\k+0.8,#2+0.62);
          \draw[vgHidden, line width=0.35pt]
            (#1+\k,#2) -- (#1+\k+0.8,#2+0.62);
          \draw[vgHidden, line width=0.35pt]
            (#1+\k+0.4,#2) -- (#1+\k+0.8,#2+0.31);
          \draw[vgHidden, line width=0.35pt]
            (#1+\k,#2+0.31) -- (#1+\k+0.4,#2+0.62);
        \fi
      \fi
    }%
  }

  % ---------------------------------------------------------------------------
  % Four same-length target-derived windows
  % ---------------------------------------------------------------------------
  \drawwin{1}{2.90}{vgKept}   % clear of hidden interval -> kept
  \drawwin{4}{1.95}{vgExcl}   % overlaps hidden interval -> excluded
  \drawwin{6}{1.00}{vgExcl}   % overlaps hidden interval -> excluded
  \drawwin{9}{0.05}{vgKept}   % clear of hidden interval -> kept

  % ---------------------------------------------------------------------------
  % Verdict labels
  % ---------------------------------------------------------------------------
  \node[anchor=west, font=\scriptsize\bfseries, text=vgKept]
    at (13.9,3.20) {kept};
  \node[anchor=west, font=\scriptsize\bfseries, text=vgExcl]
    at (13.9,2.25) {excluded};
  \node[anchor=west, font=\scriptsize\bfseries, text=vgExcl]
    at (13.9,1.30) {excluded};
  \node[anchor=west, font=\scriptsize\bfseries, text=vgKept]
    at (13.9,0.35) {kept};

\end{tikzpicture}}
  \caption{Dependency-window leakage control for tabular models.
Each tabular training row depends on a fixed target-derived window of nearby target values or summaries.
Rows whose window overlaps the hidden target interval are excluded; rows whose window stays fully outside it are kept.}
  \label{fig:loco_dependency_window}
\end{figure}

This filter applies to target-derived features, not to external predictors. Wind, SST, satellite,
and current covariates may remain available during the artificial gap because they are not the
withheld in-situ target. TS-ICL also differs from the tabular models: it is not trained on
per-gap rows, but receives the visible target sequence and time-indexed covariate channels with
the target values inside the artificial gap masked.

\subsection{Evaluation metrics, stratification, and uncertainty}
\label{sec:metrics}

The primary metric is mean absolute error (MAE), computed on the modelling scale of each target.
Chlorophyll is evaluated in $\log_{10}(\texttt{chl\_mean})$ space; back-transformed
mg~m\textsuperscript{$-3$} values are used only for physical interpretation. Dissolved oxygen is
evaluated directly in mg/L, with no log transform.

Pooled averages are not sufficient to evaluate a reconstruction method. A method may perform well
on short background gaps but fail during high-\chl{} episodes or in the tails of the oxygen
distribution. Results are therefore stratified after reconstruction. Chlorophyll errors are
examined by gap length, season, event/background status, and predictor availability. Oxygen
errors are examined by gap length, season, empirical oxygen quantile band, and tail-persistence
class. 

Uncertainty is estimated with a gap-cluster bootstrap. The resampling unit is the artificial gap,
not the individual day, and each bootstrap replicate recomputes the relevant error summary after
resampling gaps. This keeps multiple days errors from the same hidden interval grouped together and
avoids treating within-gap daily errors as independent observations.

\subsection{Covariate-signal checks}
\label{sec:covariate_signal}

The benchmark also examines whether external covariates carry useful reconstruction information.
For tabular tree ensembles, feature-importance summaries are used as descriptive diagnostics of
which engineered variables the fitted model relies on. These summaries are not interpreted as
causal attribution: correlated predictors, shared seasonality, and overlapping physical
mechanisms can all make importance scores unstable or non-unique.

For TS-ICL, covariate evidence is checked with negative-control placebo sequences. The correctly
aligned covariate sequence is compared with four disrupted versions: a wrong-lag shift, a
season-shuffled sequence, a year-shifted sequence, and a randomly permuted sequence. A covariate
family is treated as carrying useful time-aligned signal only if the correctly aligned version
beats all four placebo versions, with the gap-cluster bootstrap confidence interval for the
improvement excluding zero. These tests separate physically timed covariate information from the
generic effect of adding another input channel to the sequence model.